\documentclass[12pt]{report}
\usepackage{xcolor}

\title{\textbf{ByteRoute} \\ \large SPE Report}

\author{
  Babini Stefano\\
  0001125127 \\
  \texttt{stefano.babini5@studio.unibo.it}
}

\date{\today}

\usepackage[T1]{fontenc}
\usepackage[utf8]{inputenc}
\usepackage[english]{babel}
\usepackage{float}
\pagestyle{plain}
\usepackage{graphicx}
\usepackage{hyperref}
\usepackage{caption}
\usepackage{subcaption}
\usepackage{tikz}
\usepackage[export]{adjustbox}
\usetikzlibrary{petri,positioning}
\usepackage{listings}
\usepackage{titlesec}

\titleformat{\chapter}[display]{\normalfont\bfseries}{}{0pt}{\Large}

\hypersetup{
    colorlinks=true,
    linkcolor=black,
    filecolor=magenta,
    urlcolor=cyan,
}

\graphicspath{{images/}}

\renewcommand{\lstlistingname}{Code}
\renewcommand{\lstlistlistingname}{List of \lstlistingname s}

% Color setup
\definecolor{codegreen}{rgb}{0,0.6,0}
\definecolor{codegray}{rgb}{0.5,0.5,0.5}
\definecolor{codepurple}{rgb}{0.58,0,0.82}
\definecolor{backcolour}{rgb}{0.95,0.95,0.92}

% Code listing style
\lstdefinestyle{mystyle}{
  backgroundcolor=\color{backcolour},
  commentstyle=\color{codegreen},
  keywordstyle=\color{magenta},
  numberstyle=\tiny\color{codegray},
  stringstyle=\color{codepurple},
  basicstyle=\ttfamily\footnotesize,
  breakatwhitespace=false,
  breaklines=true,
  captionpos=b,
  keepspaces=true,
  numbers=left,
  numbersep=5pt,
  showspaces=false,
  showstringspaces=false,
  showtabs=false,
  tabsize=2
}

\lstset{style=mystyle}

\begin{document}

\maketitle
\newpage
\tableofcontents{}
\newpage

\chapter{Introduction}

This document presents a technical review of the \textbf{ByteRoute} monorepo, a multi-application project organized around network traffic analysis, visualization, and client tooling. The repository combines a TypeScript/Node.js backend, a Vue-based dashboard, a Go command-line client, and shared packages in a single coordinated workspace.

The objective of this report is to evaluate how the monorepo structure supports software quality, collaboration, and delivery in a distributed-systems-oriented project. The project is part of the SPE (Software Performance Engineering) and ASW (Architetture Software per il Web) university curricula and targets real-time visibility of internet traffic destinations.

\section{Review Scope}

The analysis focuses on:
\begin{itemize}
  \item architectural decomposition across applications and packages;
  \item development workflow and dependency management;
  \item testing, coverage, and CI/CD integration;
  \item strengths, limitations, and potential improvements.
\end{itemize}

\chapter{Repository Overview}

\section{Top-Level Organization}

The project follows a workspace-driven layout with clear separation of concerns:
\begin{itemize}
  \item \texttt{apps/backend}: Node.js/Express backend services and real-time APIs via Socket.IO;
  \item \texttt{apps/dashboard}: frontend dashboard built with Vite and Vue;
  \item \texttt{apps/client-go}: Go-based lightweight Linux client using libpcap;
  \item \texttt{packages/shared}: shared contracts, models, and reusable types;
  \item \texttt{scripts}: workspace-level automation scripts.
\end{itemize}

This structure makes ownership boundaries explicit while preserving high discoverability of shared resources.

\section{Monorepo Tooling}

The repository uses \texttt{pnpm} workspaces together with Turbo orchestration. This combination provides:
\begin{itemize}
  \item deterministic dependency installation;
  \item efficient task execution and caching;
  \item consistent build/test commands across projects.
\end{itemize}

\chapter{System Architecture}

\section{End-to-End Data Flow}

The system implements a near-real-time pipeline:
\begin{itemize}
  \item the Go client captures packets via libpcap, extracts destination IPs, deduplicates, and sends batched connections to the backend;
  \item the backend receives batches, enriches entries with GeoIP metadata, persists records in MongoDB, and broadcasts updates via Socket.IO;
  \item the Vue dashboard consumes real-time events, renders a world map of traffic flows, shows live connection lists, statistics by ASN/country, and time-based charts.
\end{itemize}

Contracts and models in the shared package reduce mismatch risks between producer (client), processor (backend), and consumer (dashboard).

\chapter{Component Analysis}

\section{Backend Application}

The backend application is implemented in TypeScript (Node.js/Express) and organized by domain-focused modules (auth, controllers, routes, services, middleware, and infrastructure). It includes:
\begin{itemize}
  \item GeoIP enrichment and Socket.IO broadcasting for live updates;
  \item MongoDB persistence for connection records and analytics;
  \item dedicated unit and feature test directories;
  \item scripts for coverage aggregation and database preparation;
  \item Docker support for containerized execution.
\end{itemize}

This reflects a mature service architecture with explicit separation between HTTP interfaces, real-time channels, business logic, and persistence concerns.

\section{Dashboard Application}

The dashboard follows a modern frontend structure with dedicated folders for components, composables, stores, services, and views. This promotes:
\begin{itemize}
  \item reusable UI logic and predictable state management;
  \item live visualization of geospatial traffic flows and trend charts;
  \item improved testability of isolated view logic.
\end{itemize}

Integration with Vite and Vitest supports fast iteration and frontend quality checks.

\section{Go Client}

The Go client introduces language diversity in the same repository while preserving modularity through \texttt{cmd/} and \texttt{internal/} packages. It is oriented to traffic capture on Linux via libpcap, destination extraction and deduplication, batching, backend communication, and runtime metrics collection.

Keeping the client close to server and shared contracts reduces integration drift and simplifies release coordination.

\section{Shared Package}

The shared package centralizes common models and types. In a multi-application environment, this is critical to avoid contract divergence and to ensure compile-time compatibility between producers and consumers.

\chapter{Quality and Delivery}

\section{Testing and Coverage}

The monorepo shows evidence of a layered testing approach:
\begin{itemize}
  \item backend unit and feature/e2e coverage;
  \item dashboard unit testing via Vitest;
  \item aggregated coverage reports surfaced through Codecov.
\end{itemize}

A unified quality baseline across applications is one of the key advantages of this monorepo setup.

\section{CI/CD Workflows}

The repository integrates GitHub Actions for automation:
\begin{itemize}
  \item release workflow builds documentation artifacts (\texttt{docs/SPE-report.pdf}, \texttt{docs/ASW-report.pdf}) prior to semantic release and publishes them as release assets;
  \item build/test workflow validates LaTeX compilation by running \texttt{make pdf} in \texttt{docs/};
  \item optional AI-assisted workflow (\texttt{.github/workflows/ai-docs-update.yml}) can update \texttt{docs/SPE-report.tex} and \texttt{docs/ASW-report.tex}, compile them, and open a pull request using GitHub Models with \texttt{GITHUB\_TOKEN}.
\end{itemize}

Selective execution based on changed paths and workspace awareness keeps pipelines efficient as the repository grows.

\section{Developer Experience}

The repository includes common governance and automation files (commit linting, semantic release configuration, renovate settings, and workspace synchronization scripts). These conventions contribute to:
\begin{itemize}
  \item consistent contribution standards;
  \item easier onboarding;
  \item safer dependency lifecycle management.
\end{itemize}

\chapter{Critical Evaluation}

\section{Strengths}

\begin{itemize}
  \item \textbf{Unified lifecycle}: build, test, documentation, and release are coordinated from one workspace.
  \item \textbf{Contract consistency}: shared models reduce mismatch across backend, dashboard, and client.
  \item \textbf{Operational readiness}: Docker, CI, coverage, and release assets indicate deployment awareness.
  \item \textbf{Real-time experience}: Socket.IO enables live telemetry to the dashboard.
  \item \textbf{Observability potential}: persisted records in MongoDB support statistics by ASN/country and time-series analysis.
\end{itemize}

\section{Potential Risks}

\begin{itemize}
  \item \textbf{Cross-language complexity}: TypeScript and Go require dual expertise and toolchain maintenance.
  \item \textbf{Packet capture overhead}: libpcap may impose CPU overhead and require elevated privileges on some systems.
  \item \textbf{GeoIP accuracy}: enrichment quality depends on database freshness and may affect analytics precision.
  \item \textbf{Pipeline coupling}: broad CI jobs may become slower as the repository scales.
  \item \textbf{Shared package impact}: breaking contract changes can affect multiple applications simultaneously.
\end{itemize}

\section{Recommended Improvements}

\begin{itemize}
  \item enforce stricter contract versioning rules for shared models and automate compatibility checks;
  \item expand architecture documentation for cross-component data flows and real-time event schemas;
  \item adopt selective CI execution based on affected projects and optimize cache usage;
  \item track non-functional metrics (build time, test duration, packet capture throughput, dashboard render latency);
  \item schedule GeoIP database refresh and validate enrichment accuracy in CI.
\end{itemize}

\chapter{Conclusion}

The ByteRoute monorepo demonstrates a well-structured and production-oriented organization. It balances modularity and integration by combining independent applications with shared contracts and centralized tooling. The explicit CI/CD around documentation and release assets strengthens reproducibility and communication.

From a university project perspective, this repository is a strong example of modern multi-service software engineering practice, with particular value in maintainability, consistency, and end-to-end coordination for network traffic monitoring.

\end{document}
